% Ch10 WS2 -- types of solids, metals, alloys, network solids. Blocks 3-4.
\documentclass{apchem-worksheet}

\apcunitword{Chapter}
\apcunit{10}
\apcunittitle{Liquids and Solids}
\apcdoctitle{Worksheet 2 \textbullet\ Solids}
\apcsubtitle{Zumdahl \S10.3--10.7 \textbullet\ name the particles and the force holding them}

\begin{document}
\wsheader[For every property question, name the \emph{particles} occupying
the lattice and the \emph{force} between them before naming the property.]

\problem[]{Classify each solid and predict whether it conducts as a solid:}
\begin{parts}
  \item \ce{W}, mp \SI{3422}{\celsius}, malleable, lustrous:
        \answerblank[5cm]{metallic; yes}
  \item \ce{SiC}, mp \SI{2730}{\celsius}, extremely hard:
        \answerblank[5cm]{network covalent; no}
  \item \ce{CaF2}, mp \SI{1418}{\celsius}, brittle:
        \answerblank[5cm]{ionic; no}
  \item \ce{S8}, mp \SI{115}{\celsius}, soft, brittle:
        \answerblank[5cm]{molecular; no}
\end{parts}

\problem[]{Explain, at the particle level, why metals are malleable but
ionic solids are brittle. \answerlines[4]{In a metal the cations are
surrounded by a delocalized electron sea. When layers slide, the sea flows
with them and no repulsion develops, so the metal deforms. In an ionic solid
a layer shift brings like charges into alignment --- cation above cation ---
and the sudden repulsion splits the crystal along a clean plane.}}

\problem[]{Diamond and graphite are both pure carbon.}
\begin{parts}
  \item Hybridization of carbon in each:
        \answerblank[5.3cm]{diamond sp$^3$; graphite sp$^2$}
  \item Which conducts electricity, and why?
        \answerlines[2]{Graphite --- each sp$^2$ carbon has a leftover $p$
        electron, and these are delocalized across the layer, giving mobile
        charge carriers. Diamond's electrons are all localized in
        $\sigma$ bonds.}
  \item Why is graphite soft enough to use as a lubricant while diamond is
        the hardest known material? \answerlines[3]{Graphite's layers are
        strongly bonded internally but held to one another only by weak
        dispersion forces, so they slide easily. Diamond is a
        three-dimensional covalent network, so deforming it requires
        breaking covalent bonds in every direction.}
\end{parts}

\problem[]{Alloys:}
\begin{parts}
  \item Brass substitutes zinc for copper; steel inserts carbon into iron.
        Name each type. \answerblank[7.4cm]{brass substitutional; steel
        interstitial}
  \item Why is steel harder than pure iron?
        \answerlines[2]{Interstitial carbon atoms occupy holes between iron
        atoms and block the layers from sliding past one another.}
  \item Does alloying destroy electrical conductivity? Explain.
        \answerlines[2]{No --- the delocalized electron sea remains intact,
        so charge carriers are still mobile.}
\end{parts}

\problem[]{Molecular and ionic solids:}
\begin{parts}
  \item Melting ice breaks which forces? Support with numbers.
        \answerlines[2]{Hydrogen bonds between molecules
        ($\approx$\SI{20}{\kilo\joule\per\mole}), not the O--H covalent
        bonds ($\approx$\SI{460}{\kilo\joule\per\mole}) inside them.}
  \item Explain why ice is less dense than liquid water.
        \answerlines[2]{Hydrogen bonding locks the molecules into an open
        hexagonal lattice containing more empty space than the liquid
        arrangement.}
  \item Why does \ce{NaCl} conduct when molten but not as a solid?
        \answerlines[2]{Conduction requires mobile charge carriers. In the
        solid the ions are locked in the lattice; melting frees them to
        move.}
\end{parts}

\problem[]{Rank by increasing melting point and justify each step:
\ce{Ar}, \ce{HCl}, \ce{KCl}, \ce{diamond}.
\answerlines[4]{\ce{Ar} $<$ \ce{HCl} $<$ \ce{KCl} $<$ diamond. Argon is
molecular with dispersion only; \ce{HCl} is molecular with dipole--dipole as
well; \ce{KCl} is ionic with a full lattice of electrostatic attractions;
diamond is a covalent network requiring covalent bonds to break.}}

\begin{spiralstrip}{Chapter 10 \textbullet\ blocks 1--2}
  \item Strongest IMF in \ce{HF}: \answerblank[3.8cm]{hydrogen bonding}
  \item Water in glass: concave or convex?
        \answerblank[2.2cm]{concave}
  \item Higher vapour pressure means IMFs are:
        \answerblank[2cm]{weaker}
  \item Which has stronger dispersion, \ce{I2} or \ce{F2}?
        \answerblank[2cm]{\ce{I2}}
  \item Surface tension increases with what?
        \answerblank[3.4cm]{stronger IMFs}
\end{spiralstrip}

\end{document}
